\documentclass[11pt,a4paper]{article}

\usepackage[utf8]{inputenc}
\usepackage[T1]{fontenc}
\usepackage[margin=2.5cm]{geometry}
\usepackage{amsmath}
\usepackage{graphicx}
\usepackage{booktabs}
\usepackage{xcolor}
\usepackage{url}

% Visible placeholder: remove every \todo before submitting.
\newcommand{\todo}[1]{\textcolor{red}{[TODO: #1]}}

\title{Tiered Memory Management Beyond Hotness\\
\large CS-471: Project Status Report}

\author{
Mayeul Besnard\\
EPFL\\
\texttt{mayeul.besnard@epfl.ch}
\and
Sergio Boffi\\
EPFL\\
\texttt{sergio.boffi@epfl.ch}
}

\date{November 12, 2026}

\begin{document}

\maketitle

\tableofcontents
\newpage

% =====================================================================
\section{Research Problem}
Tiered memory systems combine a small, fast tier (local DRAM) with a larger,
cheaper but slower tier (e.g., CXL-attached memory), so the way data is
distributed across the two tiers determines how many memory stalls an
application suffers. State-of-the-art tiering systems (TPP, NBT, Nomad and
Colloid) make this decision based on access frequency, or \emph{hotness}.

Liu et al.\ (OSDI 2025) argue that hot data is not necessarily
performance-critical. Out-of-order cores overlap independent misses, so accesses
with high memory-level parallelism (MLP) hide most of the slow-tier latency,
while low-MLP accesses, such as pointer chasing, pay all of it. The authors
introduce the Amortized Offcore Latency, $\text{AOL} = \text{Latency}/\text{MLP}$,
model the slowdown as $S = P \times K(\text{AOL})$, and build on it SOAR, an
AOL-guided first-placement policy.

Our project aims to reproduce these claims by answering four questions:
\begin{itemize}
  \item \textbf{Q1:} Can hotness be a faithful metric to measure performance?
  \item \textbf{Q2:} If not, how can a better metric be defined?
  \item \textbf{Q3:} With this metric, can we improve first data placement (SOAR)?
  \item \textbf{Q4:} If time permits, can these elements improve PolyRack?
\end{itemize}

The main focus of the project has not changed since the proposal. What has
changed is our understanding of the platform: the server available to us is
based on Intel Haswell, an older generation than the Skylake and Sapphire Rapids
machines used by the authors, which affects how the metrics can be measured
(Section~2.1).

% =====================================================================
\section{Questions, Hypotheses, and Tests}

\subsection{Methodology and Testing Environment Updates}
\todo{Machine actually used and how it differs from the assumptions of the
proposal and from the paper.}

\todo{Changes to the methodology: counters used, measurement procedure, how the
slow tier is emulated.}

\subsection{Current Status of Research Questions}

\subsubsection{Q1: Can hotness be a faithful metric to measure performance?}
\textbf{Insight:} \todo{I1 from the proposal.}

\textbf{Experiment:} \todo{E1, and the test of H1$'$.}

\textbf{Status:} \todo{Done, partial or not started.}

\subsubsection{Q2: If not, how can a better metric be defined?}
\textbf{Insight:} \todo{I2 from the proposal.}

\textbf{Experiment:} \todo{E2, and the evaluation of H2$'$.}

\textbf{Status:} \todo{Done, partial or not started.}

\subsubsection{Q3: With this metric, can we improve first data placement (SOAR)?}
\textbf{Insight:} \todo{I3 from the proposal.}

\textbf{Experiment:} \todo{E3.}

\textbf{Status:} \todo{Done, partial or not started.}

\subsubsection{Q4: Can these elements improve PolyRack?}
\textbf{Insight:} \todo{I4 from the proposal.}

\textbf{Experiment:} \todo{E4.}

\textbf{Status:} \todo{Done, partial, not started, or dropped.}

\subsection{Anomalies and Updates to the Hypothesis Tree}
\todo{Any anomaly observed, and whether the hypothesis tree changes.
Include the updated figure if it does.}

% =====================================================================
\section{Progress, Results, and Analysis}

\subsection{Experimental Setup}
\todo{What has been brought up and verified on the machine.}

\subsection{E1: Hotness versus Criticality}
\todo{Configurations, number of runs, results and conclusion on H1.}

\begin{figure}[h]
\centering
\fbox{\parbox[c][4cm][c]{0.7\linewidth}{\centering \todo{E1 plot}}}
\caption{\todo{Caption.}}
\label{fig:e1}
\end{figure}

\subsection{E2: Slowdown Predictors (Preliminary)}
\todo{Calibration of $a,b$ and comparison of the predictors.}

\begin{table}[h]
\centering
\begin{tabular}{lcc}
\toprule
Predictor & Pearson & MAE \\
\midrule
$P$ & \todo{--} & \todo{--} \\
$P \times K$ & \todo{--} & \todo{--} \\
Stalls per miss (H2$'$) & \todo{--} & \todo{--} \\
Hotness & \todo{--} & \todo{--} \\
\bottomrule
\end{tabular}
\caption{\todo{Caption.}}
\label{tab:e2}
\end{table}

\subsection{SOAR: Feasibility}
\todo{What the authors' artifact requires, what is needed to run it on our
machine, and any first result.}

% =====================================================================
\section{Challenges}

\subsection{Hardware Differences}
\todo{Challenges coming from the CPU and its counters.}

\subsection{Tiered Memory Emulation and Baselines}
\todo{Challenges with the slow tier, memory tiers and custom kernels.}

\subsection{Authors' Artifact}
\todo{Challenges in porting the code.}

\subsection{Help Needed from the Teaching Team}
\todo{Concrete requests.}

% =====================================================================
\section{Plan}
\todo{Say whether the project is on schedule with respect to the proposal.}

\begin{table}[h]
\centering
\begin{tabular}{lll}
\toprule
Date & Milestone & Status \\
\midrule
\todo{date} & \todo{E2 completed} & \todo{--} \\
\todo{date} & \todo{SOAR results} & \todo{--} \\
Nov.~27 & \todo{E3 on selected workloads and ratios} & \todo{--} \\
Dec.~4 & \todo{(if time permits) E4 on PolyRack} & \todo{--} \\
Dec.~10 & Poster deadline & \\
Dec.~17 & Poster session & \\
Dec.~20 & Final report deadline & \\
\bottomrule
\end{tabular}
\caption{Updated milestones.}
\label{tab:plan}
\end{table}

\todo{Division of the work between the team members.}

\end{document}
